\subsection{Support Vector Machine}
\label{sec:support_vector_machine}

\begin{table*}[htbp]
\caption{Hyperparameter Search Space for Support Vector Machine.}
\label{tab:VI}
\centering
\begin{tabular}{llc}
\toprule
Component / Hyperparameter & Evaluated Values & Count \\
\midrule
Feature Scaler & None, MinMaxScaler & 2 \\
Dimensionality Reduction (PCA) & None, 0.50, 0.75 & 3 \\
Regularization Parameter & 0.01, 0.1, 1, 10 & 4 \\
Kernel Function & 'linear', 'rbf', 'poly' & 3 \\
Polynomial Degree & 2, 3 & 2 \\
Kernel Coefficient & 'scale', 'auto' & 2 \\
\midrule
\textbf{Total Grid Combinations} & -- & \textbf{288 (1,440 fits)} \\
\bottomrule
\end{tabular}
\end{table*}

Pengklasifikasi SVM mengonstruksikan bidang pemisah berjarak terjauh (maximum margin hyperplane) untuk memisahkan enam kelas interseksional pada ruang representasi laten berdimensi tinggi. Fungsi kernel memproyeksikan distribusi data nonlinier dari fusi multi-domain ke ruang fitur berdimensi lebih tinggi, mencakup opsi kernel linear, Radial Basis Function (RBF), dan polinomial. Formulasi kernel polinomial derajat dua pada \eqref{eq:8} menghitung produk titik $\langle \mathbf{x}_i, \mathbf{x}_j \rangle$ antara pasangan vektor fitur $\mathbf{x}_i$ dan $\mathbf{x}_j$, di mana $\gamma$ menyatakan parameter penskalaan kernel, $d = 2$ menentukan derajat pemetaan, dan konstanta intercept $\text{coef0} = 0.0$ mengikuti nilai bawaan scikit-learn:
\begin{equation}
K(\mathbf{x}_i, \mathbf{x}_j) = (\gamma \langle \mathbf{x}_i, \mathbf{x}_j \rangle + \text{coef0})^d, \quad d = 2
\label{eq:8}
\end{equation}
Pemetaan kuadratik ini memfasilitasi penangkapan interaksi nonlinier antardimensi visual secara efisien.

Optimasi hyperparameter menelusuri rentang parameter penalti $C \in \{0.01, 0.1, 1, 10\}$, ragam fungsi kernel, koefisien $\gamma$, derajat polinomial, serta opsi penskalaan dan PCA, menghasilkan 288 kombinasi konfigurasi atau 1,440 total fits pada prosedur 5-Fold Stratified Cross-Validation sesuai Table~\ref{tab:VI}.

